\documentclass[10pt,a4paper]{amsart}
\usepackage[margin=19mm]{geometry}
\usepackage{amsmath,amssymb,mathtools}
\usepackage[scr=rsfs,cal=euler]{mathalfa}
\usepackage{xfrac}
\usepackage[unicode,hidelinks,bookmarks=false]{hyperref}
\newcommand{\ie}{i.e.\ }
\newcommand{\cf}{cf.\ }
\newcommand{\resp}{respectively\ }
\newcommand{\Subsection}{\subsection}
\providecommand{\nobreakdash}{\nobreak\hbox{-}}
\setlength{\emergencystretch}{2em}
\multlinegap0pt
\usepackage{bm}
\usepackage{bbm}
\usepackage{dsfont}
\usepackage{xspace}
\usepackage{cancel}
\usepackage{mathtools}
\usepackage{amsthm}
\makeatletter
\@ifundefined{theorem}{\newtheorem{theorem}{Theorem}}{}
\@ifundefined{lemma}{\newtheorem{lemma}[theorem]{Lemma}}{}
\@ifundefined{proposition}{\newtheorem{proposition}[theorem]{Proposition}}{}
\makeatother
\newcommand{\azero}{\mathbf{0}}
\title[The Algebraic Boundary of Graph Elliptopes]{The Algebraic Boundary of Graph Elliptopes}
\author{Monique Laurent, Francesco Maria Mascarin, Simon Telen}
\date{Source: arXiv:2605.02484v1 (CC BY 4.0)}
\begin{document}
\maketitle

\noindent\emph{Source and licence.} The Algebraic Boundary of Graph Elliptopes. Version arXiv:2605.02484v1 (CC BY 4.0), distributed under the Creative Commons Attribution 4.0 International licence, as recorded in the arXiv licence metadata for this exact version and shown on the abstract page https://arxiv.org/abs/2605.02484v1. Full rights notice: licenses/arxiv-2605.02484-CC-BY-4.0.txt.\par\smallskip

\noindent\textit{Editorial scope.} Counted unit: the Proposition whose source label is \texttt{prop:Resh} in the article (source0.tex lines 807--810, source label \texttt{prop:Resh}), delivered together with the article's own complete original proof of it (source0.tex lines 812--848, one \texttt{proof} environment of 37 lines). No mathematics is written, repaired or re-derived: statement and proof are the article's own text line for line. Required context retained uncounted: eq:hn (lines 531--533); eq:zepsilon (lines 603--605); eq:hnz (lines 610--612); prop:univleadingterm (lines 616--627); lem:z (lines 725--731); eq:Poisson (lines 787--789). Every definition, display and label of those lines is retained unchanged. Omitted from this package and not redistributed: 1--530, 534--602, 606--609, 613--615, 628--724, 732--786, 790--806, 811--811, 849--1486. Nothing inside a retained excerpt is omitted or altered: every retained body is delivered exactly as the article's lines are written, so no omission receipt of any kind accompanies this package. Citations: the retained excerpts carry no citation command at all, so no bibliography entry of the article is transcribed and no reference is redistributed. Reference adaptation: none was needed and none was made; every reference inside the delivered unit resolves inside it, so no unresolved reference or citation is produced. Printed slips kept literal: no printed word, symbol, hypothesis or number of the counted statement or of its proof was altered. Layout and class: the article's own class is \texttt{article} with packages including geometry, inputenc, lmodern, amssymb, babel, graphicx, epsfig, epstopdf, amsfonts, amsmath, braket, amsthm, chngpage, xcolor; the wrapper is the lane's pinned \texttt{amsart} editorial wrapper at 10pt on A4, which restates the theorem-like environments the retained text needs and adds the article's own macro definitions that the retained text needs (\texttt{\textbackslash azero}), with extra wrapper packages (bm, bbm, dsfont, xspace, cancel, mathtools). The delivered numbering is the unit's own. Assets: no retained span contains a figure environment, a graphics-inclusion command, a PDF-page inclusion, a table or a TikZ picture; no image file is copied into this package, the package declares no asset dependency and no image file is redistributed with it. Licence: version arXiv:2605.02484v1 of this article is distributed under the Creative Commons Attribution 4.0 International licence, as recorded in the arXiv licence metadata for this exact version and shown on the abstract page https://arxiv.org/abs/2605.02484v1. A full rights notice is delivered at licenses/arxiv-2605.02484-CC-BY-4.0.txt. Characters: every retained excerpt is pure ASCII, so the delivered engine, XeLaTeX with its default Latin Modern text font, reports no missing character.
\begin{equation}\label{eq:hn}
h_n(x) \,  = \, \prod_{\xi \in V(J_x)} g_n(\xi) \, = \, \prod_{\epsilon \in \{-1,1\}^{E_n}} \Big[ \prod_{e \in E_n} (x_e + \epsilon_e \sqrt{x_e^2-1} ) - 1 \Big].
\end{equation}

\begin{align}\label{eq:zepsilon}
z_\epsilon= \prod_{e\in E_n} \left( x_e + \epsilon_e \sqrt{x_e^2-1} \right)^{|\epsilon_e|} \ \text{ for any } \epsilon \in \{\pm1, 0\}^{E_n}.
\end{align}

\begin{align}\label{eq:hnz}
h_n=\prod_{\epsilon\in \{\pm 1\}^{E_n}} (z_\epsilon -1).
\end{align}

\begin{proposition} \label{prop:univleadingterm}
 Consider two distinct edges $e_0,e_1\in E_n$. When viewed as a univariate polynomial in $x_{e_0}$, the polynomial $h_n$ 
has degree $2^{n-1}$ with leading coefficient $2^{2^{n-1}}$, and its roots are given by the following algebraic functions of $x_e$ (for $e \in E_n \setminus \{e_0\}$):
\begin{align}\label{eq:rdelta}
 r_\delta \,  = \, \frac{z_{0,\delta} + z_{0,-\delta}}{2}  \quad \text{ for } \delta \in \{-1,1\}^{E_n \setminus \{e_0\}}, \end{align}
where $z_{0,\delta} \, = \, \prod_{e \in E \setminus \{e_0\}} (x_{e} + \delta_{e} \, \sqrt{x_{e}^2-1} )$ as in (\ref{eq:zepsilon}), with the first zero entry in the subscript corresponding to the variable $x_{e_0}$.
  Moreover,  we have
 \begin{align}\label{eq:hnroot}
 h_n= 2^{2^{n-1}} \prod_{\delta\in \{\pm 1\}^{E_n\setminus \{e_0\}}} z_{0,\delta} (x_{e_0}-r_\delta)
 =2^{2^{n-1}} \prod_{\delta\in \{\pm 1\}^{E_n\setminus \{e_0\}}: \delta_{e_1}=1}  (x_{e_0}-r_\delta)^2.
     \end{align}
    \end{proposition}

\begin{lemma}
\label{lem:z}
For $\delta\in\{\pm 1\}^{E_n\setminus \{e_0\}}$ and $\epsilon_0\in \{\pm 1\}$, there exists $\eta_{\delta}\in\{\pm 1\}$ such that 
\[
r_\delta +\epsilon_0\sqrt{r_\delta^2 -1}= z_{0,\epsilon_0\cdot \eta_\delta \cdot  \delta}.
\]
\end{lemma}

\begin{equation} \label{eq:Poisson}
{\rm Res}_y(f,g) \, = \, a_{d_f}^{d_g} \, \prod_{i = 1}^{d_f} g(r_i).
\end{equation}

\begin{proposition} \label{prop:Resh}
With the above notation, consider the polynomials $h_n \in \mathbb{C}[x_e  :  e \in E_n]$, $h_{C_p} \in \mathbb{C}[x_e : e \in F_p\cup\{e_0\}]$, and $h_{C'_q} \in \mathbb{C}[x_e  :  e \in F_q'\cup\{e_0\}]$,  where $E_n=F_p\cup F'_q$. 
We have the identity: ${\rm Res}_{x_{e_0}}(h_{C_p}, h_{C'_q}) = (2^{2^{n-1}} \, h_n)^2$. 
\end{proposition}

\begin{proof}
We compute the resultant using relation (\ref{eq:Poisson}) applied to $f= h_{C_p}$ and $g=h_{C'_q}$, viewed as univariate polynomials in $x_{e_0}$.
By Proposition \ref{prop:univleadingterm}, the roots of $h_{C_p}$ 
are given by 
$r_\delta={1\over 2}(z_{0,\delta,\azero}+z_{0,-\delta,\azero})$ for $ \delta \in \{-1,1\}^{F_p}.$
Here, the subscript of $z$ uses the   ordering of $\{e_0\} \cup E_n$ as $e_0, F_p, F_q'$. 
By Poisson's formula \eqref{eq:Poisson}, the resultant is given by the following~expression:
\[
{\rm Res}_{x_{e_0}}(h_{C_p}, h_{C_q'})= (2^{2^{p-1}})^{2^{q-1}} \, \prod_{\delta \in \{-1,1\}^{F_p}}  h_{C'_q}(r_\delta)=
2^{2^n} \prod_{\delta \in \{-1,1\}^{F_p}}  h_{C'_q}(r_\delta),  \]
where we used $p + q - 2 = n$. Using the definition (\ref{eq:hnz}) for the polynomial $h_{C_q'}$, we get
\begin{align*}
h_{C_q'}= \prod_{\delta'\in\{\pm 1\}^{F_q'}} \prod_{\epsilon_0\in \{\pm 1\}} (z_{\epsilon_0,\azero,\delta'}-1) = 
 \prod_{\delta'\in\{\pm 1\}^{F_q'}} \prod_{\epsilon_0\in \{\pm 1\}} (z_{\epsilon_0,\azero,\azero}\cdot z_{0,\azero,\delta'}-1).
 \end{align*}
By Lemma \ref{lem:z}, evaluating $z_{\epsilon_0,\azero,\azero}$ at $x_{e_0}=r_\delta$ gives:
\[
(z_{\epsilon_0,\azero,\azero})_{|x_{e_0} = r_\delta} \,=\, r_{\delta}+\epsilon_0\sqrt{r_\delta^2-1} \,= \,z_{0,\epsilon_0\cdot \eta_\delta\cdot  \delta,\azero}\ \text{   for some } \eta_\delta\in\{\pm 1\}.
\]
Then, viewing $h_{C_q'}$ as a polynomial in $x_{e_0}$ and evaluating it at $x_{e_0}=r_\delta$, 
   we obtain 
\begin{align*}
\prod_{\delta \in \{-1,1\}^{F_p}}  h_{C'_q}(r_\delta)& =
\prod_{\delta \in \{-1,1\}^{F_p}}  \prod_{\delta'\in\{\pm 1\}^{F_q'}}\prod_{\epsilon_0\in \{\pm 1\}}
(z_{0,\epsilon_0\cdot \eta_\delta\cdot  \delta,\azero}\cdot z_{0,\azero,\delta'}-1)\\
&=
\prod_{\delta \in \{-1,1\}^{F_p}}  \prod_{\delta'\in\{\pm 1\}^{F_q'}}
(z_{0,\delta,\delta'}-1)^2\\
&= \prod_{\epsilon\in \{\pm 1\}^{E_n}} (z_{0,\epsilon}-1)^2\\
&= h_n^2.
\end{align*}
Here,  we use  the identity
$\prod_\delta\prod_{\epsilon_0}(z_{0,\epsilon_0\cdot\eta_\delta \cdot \delta,\azero}\cdot z_{0,\azero,\delta'}-1)=\prod_\delta 
(z_{0,\delta,\azero}\cdot z_{0,\azero,\delta'}-1)^2$ for the second equality, combined with $z_{0,\delta,\azero}\cdot z_{0,\azero,\delta'}=z_{0,\delta,\delta'}$. 
The last equality follows using relation (\ref{eq:hn}) and the fact that $E_n=F_p\cup F_q'$.
So, we have proved that ${\rm Res}_{x_{e_0}}(h_{C_p}, h_{C_q'})=2^{2^{n}} h_n^2$, as desired.
\end{proof}

\end{document}
